\documentclass[11pt]{article}
% =============================================================================
%  Shared article preamble — tutorials, homework, solutions.
%  Same fonts, palette and notation as the Beamer decks (see preamble.tex).
%      \documentclass[11pt]{article}
%      % =============================================================================
%  Shared article preamble — tutorials, homework, solutions.
%  Same fonts, palette and notation as the Beamer decks (see preamble.tex).
%      \documentclass[11pt]{article}
%      % =============================================================================
%  Shared article preamble — tutorials, homework, solutions.
%  Same fonts, palette and notation as the Beamer decks (see preamble.tex).
%      \documentclass[11pt]{article}
%      \input{../../../style/preamble_article}
% =============================================================================
\usepackage[margin=1in]{geometry}
\usepackage{helvet}
\usepackage[default]{lato}
\usepackage{mathpazo}
\usepackage[utf8]{inputenc}
\usepackage{amsmath,amssymb,amsthm}
\usepackage{mathtools}
\usepackage{graphicx}
\usepackage{booktabs}
\usepackage{multirow}
\usepackage{enumitem}
\usepackage{xcolor}
\usepackage{tcolorbox}
\tcbuselibrary{skins,breakable}
\usepackage{listings}
\usepackage[ruled,vlined]{algorithm2e}
\usepackage{hyperref}
\usepackage{fancyhdr}
\usepackage{titlesec}

\definecolor{blue}{RGB}{0,114,178}
\definecolor{red}{RGB}{213,94,0}
\definecolor{yellow}{RGB}{240,228,66}
\definecolor{green}{RGB}{0,158,115}
\hypersetup{colorlinks=true, linkcolor=blue, urlcolor=blue, citecolor=blue}

\titleformat{\section}{\Large\bfseries\color{blue}}{\thesection}{0.6em}{}
\titleformat{\subsection}{\large\bfseries}{\thesubsection}{0.6em}{}
\pagestyle{fancy}\fancyhf{}
\rhead{\footnotesize ENEI -- INEI \,\textperiodcentered\, PEU-CD 2026}
\lfoot{\footnotesize Machine Learning I \& II \,\textperiodcentered\, Alexander Quispe}
\rfoot{\footnotesize\thepage}
\renewcommand{\headrulewidth}{0.3pt}

\newtcolorbox{derivation}[1]{enhanced, breakable, colback=blue!3, colframe=blue, fonttitle=\bfseries,
  title={#1}, boxrule=0.6pt, left=6pt, right=6pt, top=4pt, bottom=4pt, arc=2pt}
\newtcolorbox{claim}[1]{enhanced, breakable, colback=white, colframe=black, fonttitle=\bfseries,
  title={#1}, boxrule=0.6pt, left=6pt, right=6pt, top=4pt, bottom=4pt, arc=2pt}
\newtcolorbox{keypoint}{enhanced, breakable, colback=green!6, colframe=green, boxrule=0.6pt,
  left=6pt, right=6pt, top=4pt, bottom=4pt, arc=2pt}
\newtcolorbox{caution}{enhanced, breakable, colback=red!5, colframe=red, boxrule=0.6pt,
  left=6pt, right=6pt, top=4pt, bottom=4pt, arc=2pt}
\newtcolorbox{task}[1]{enhanced, breakable, colback=yellow!12, colframe=yellow!60!black, fonttitle=\bfseries,
  title={#1}, boxrule=0.6pt, left=6pt, right=6pt, top=4pt, bottom=4pt, arc=2pt}

\lstset{
  basicstyle=\ttfamily\small, keywordstyle=\color{blue}\bfseries,
  commentstyle=\color{green}, stringstyle=\color{red},
  showstringspaces=false, breaklines=true, frame=single, rulecolor=\color{black!30},
  numbers=none, xleftmargin=6pt, framexleftmargin=6pt, language=Python
}
\newtheorem{theorem}{Theorem}
\newtheorem{lemma}{Lemma}
\newtheorem{proposition}{Proposition}
\theoremstyle{definition}
\newtheorem{definition}{Definition}
\newtheorem{exercise}{Exercise}

\input{../../../style/macros}

% =============================================================================
\usepackage[margin=1in]{geometry}
\usepackage{helvet}
\usepackage[default]{lato}
\usepackage{mathpazo}
\usepackage[utf8]{inputenc}
\usepackage{amsmath,amssymb,amsthm}
\usepackage{mathtools}
\usepackage{graphicx}
\usepackage{booktabs}
\usepackage{multirow}
\usepackage{enumitem}
\usepackage{xcolor}
\usepackage{tcolorbox}
\tcbuselibrary{skins,breakable}
\usepackage{listings}
\usepackage[ruled,vlined]{algorithm2e}
\usepackage{hyperref}
\usepackage{fancyhdr}
\usepackage{titlesec}

\definecolor{blue}{RGB}{0,114,178}
\definecolor{red}{RGB}{213,94,0}
\definecolor{yellow}{RGB}{240,228,66}
\definecolor{green}{RGB}{0,158,115}
\hypersetup{colorlinks=true, linkcolor=blue, urlcolor=blue, citecolor=blue}

\titleformat{\section}{\Large\bfseries\color{blue}}{\thesection}{0.6em}{}
\titleformat{\subsection}{\large\bfseries}{\thesubsection}{0.6em}{}
\pagestyle{fancy}\fancyhf{}
\rhead{\footnotesize ENEI -- INEI \,\textperiodcentered\, PEU-CD 2026}
\lfoot{\footnotesize Machine Learning I \& II \,\textperiodcentered\, Alexander Quispe}
\rfoot{\footnotesize\thepage}
\renewcommand{\headrulewidth}{0.3pt}

\newtcolorbox{derivation}[1]{enhanced, breakable, colback=blue!3, colframe=blue, fonttitle=\bfseries,
  title={#1}, boxrule=0.6pt, left=6pt, right=6pt, top=4pt, bottom=4pt, arc=2pt}
\newtcolorbox{claim}[1]{enhanced, breakable, colback=white, colframe=black, fonttitle=\bfseries,
  title={#1}, boxrule=0.6pt, left=6pt, right=6pt, top=4pt, bottom=4pt, arc=2pt}
\newtcolorbox{keypoint}{enhanced, breakable, colback=green!6, colframe=green, boxrule=0.6pt,
  left=6pt, right=6pt, top=4pt, bottom=4pt, arc=2pt}
\newtcolorbox{caution}{enhanced, breakable, colback=red!5, colframe=red, boxrule=0.6pt,
  left=6pt, right=6pt, top=4pt, bottom=4pt, arc=2pt}
\newtcolorbox{task}[1]{enhanced, breakable, colback=yellow!12, colframe=yellow!60!black, fonttitle=\bfseries,
  title={#1}, boxrule=0.6pt, left=6pt, right=6pt, top=4pt, bottom=4pt, arc=2pt}

\lstset{
  basicstyle=\ttfamily\small, keywordstyle=\color{blue}\bfseries,
  commentstyle=\color{green}, stringstyle=\color{red},
  showstringspaces=false, breaklines=true, frame=single, rulecolor=\color{black!30},
  numbers=none, xleftmargin=6pt, framexleftmargin=6pt, language=Python
}
\newtheorem{theorem}{Theorem}
\newtheorem{lemma}{Lemma}
\newtheorem{proposition}{Proposition}
\theoremstyle{definition}
\newtheorem{definition}{Definition}
\newtheorem{exercise}{Exercise}

% =============================================================================
%  Notation — one convention for all lectures
%
%  scalars      x, y, w           plain italic
%  vectors      \vx \vy \vw       bold lower-case
%  matrices     \vX \vW \vSigma   bold upper-case
%  parameters   \vbeta \vtheta    bold Greek
%  transpose    \vX\T             \top
% =============================================================================
\newcommand{\R}{\mathbb{R}}
\newcommand{\E}{\mathbb{E}}
\newcommand{\Prob}{\mathbb{P}}
\DeclareMathOperator{\Var}{Var}
\DeclareMathOperator{\Cov}{Cov}
\DeclareMathOperator{\tr}{tr}
\DeclareMathOperator{\diag}{diag}
\DeclareMathOperator{\rank}{rank}
\DeclareMathOperator{\sign}{sign}
\DeclareMathOperator{\col}{col}
\DeclareMathOperator{\RSS}{RSS}
\DeclareMathOperator*{\argmin}{arg\,min}
\DeclareMathOperator*{\argmax}{arg\,max}
\newcommand{\norm}[1]{\left\lVert #1 \right\rVert}
\newcommand{\abs}[1]{\left\lvert #1 \right\rvert}
\newcommand{\inner}[2]{\left\langle #1,\, #2 \right\rangle}
\newcommand{\ind}[1]{\mathbf{1}\!\left\{#1\right\}}
\newcommand{\T}{^{\top}}
\newcommand{\inv}{^{-1}}
\newcommand{\half}{\tfrac{1}{2}}
\newcommand{\eps}{\varepsilon}
\newcommand{\Dset}{\mathcal{D}}
\newcommand{\Hset}{\mathcal{H}}
\newcommand{\Lcal}{\mathcal{L}}
\newcommand{\Ncal}{\mathcal{N}}

% bold vectors
\newcommand{\vx}{\mathbf{x}}   \newcommand{\vy}{\mathbf{y}}   \newcommand{\vz}{\mathbf{z}}
\newcommand{\vw}{\mathbf{w}}   \newcommand{\vu}{\mathbf{u}}   \newcommand{\vv}{\mathbf{v}}
\newcommand{\vb}{\mathbf{b}}   \newcommand{\vh}{\mathbf{h}}   \newcommand{\vr}{\mathbf{r}}
\newcommand{\vc}{\mathbf{c}}   \newcommand{\vf}{\mathbf{f}}   \newcommand{\vg}{\mathbf{g}}
\newcommand{\vp}{\mathbf{p}}   \newcommand{\vq}{\mathbf{q}}   \newcommand{\vs}{\mathbf{s}}
\newcommand{\va}{\mathbf{a}}   \newcommand{\vd}{\mathbf{d}}   \newcommand{\vt}{\mathbf{t}}
\newcommand{\vk}{\mathbf{k}}   \newcommand{\vm}{\mathbf{m}}   \newcommand{\vn}{\mathbf{n}}
\newcommand{\vzero}{\mathbf{0}} \newcommand{\vone}{\mathbf{1}}
% bold matrices
\newcommand{\vX}{\mathbf{X}}   \newcommand{\vY}{\mathbf{Y}}   \newcommand{\vZ}{\mathbf{Z}}
\newcommand{\vW}{\mathbf{W}}   \newcommand{\vH}{\mathbf{H}}   \newcommand{\vI}{\mathbf{I}}
\newcommand{\vA}{\mathbf{A}}   \newcommand{\vB}{\mathbf{B}}   \newcommand{\vC}{\mathbf{C}}
\newcommand{\vD}{\mathbf{D}}   \newcommand{\vP}{\mathbf{P}}   \newcommand{\vQ}{\mathbf{Q}}
\newcommand{\vU}{\mathbf{U}}   \newcommand{\vV}{\mathbf{V}}   \newcommand{\vS}{\mathbf{S}}
\newcommand{\vM}{\mathbf{M}}   \newcommand{\vK}{\mathbf{K}}   \newcommand{\vG}{\mathbf{G}}
% bold Greek
\newcommand{\vbeta}{\boldsymbol{\beta}}     \newcommand{\valpha}{\boldsymbol{\alpha}}
\newcommand{\vtheta}{\boldsymbol{\theta}}   \newcommand{\vmu}{\boldsymbol{\mu}}
\newcommand{\vSigma}{\boldsymbol{\Sigma}}   \newcommand{\vLambda}{\boldsymbol{\Lambda}}
\newcommand{\veps}{\boldsymbol{\varepsilon}} \newcommand{\vphi}{\boldsymbol{\phi}}
\newcommand{\vgamma}{\boldsymbol{\gamma}}   \newcommand{\vlambda}{\boldsymbol{\lambda}}
\newcommand{\vdelta}{\boldsymbol{\delta}}   \newcommand{\vnu}{\boldsymbol{\nu}}
\newcommand{\vxi}{\boldsymbol{\xi}}         \newcommand{\vpi}{\boldsymbol{\pi}}
\newcommand{\vomega}{\boldsymbol{\omega}}



% =============================================================================
\usepackage[margin=1in]{geometry}
\usepackage{helvet}
\usepackage[default]{lato}
\usepackage{mathpazo}
\usepackage[utf8]{inputenc}
\usepackage{amsmath,amssymb,amsthm}
\usepackage{mathtools}
\usepackage{graphicx}
\usepackage{booktabs}
\usepackage{multirow}
\usepackage{enumitem}
\usepackage{xcolor}
\usepackage{tcolorbox}
\tcbuselibrary{skins,breakable}
\usepackage{listings}
\usepackage[ruled,vlined]{algorithm2e}
\usepackage{hyperref}
\usepackage{fancyhdr}
\usepackage{titlesec}

\definecolor{blue}{RGB}{0,114,178}
\definecolor{red}{RGB}{213,94,0}
\definecolor{yellow}{RGB}{240,228,66}
\definecolor{green}{RGB}{0,158,115}
\hypersetup{colorlinks=true, linkcolor=blue, urlcolor=blue, citecolor=blue}

\titleformat{\section}{\Large\bfseries\color{blue}}{\thesection}{0.6em}{}
\titleformat{\subsection}{\large\bfseries}{\thesubsection}{0.6em}{}
\pagestyle{fancy}\fancyhf{}
\rhead{\footnotesize ENEI -- INEI \,\textperiodcentered\, PEU-CD 2026}
\lfoot{\footnotesize Machine Learning I \& II \,\textperiodcentered\, Alexander Quispe}
\rfoot{\footnotesize\thepage}
\renewcommand{\headrulewidth}{0.3pt}

\newtcolorbox{derivation}[1]{enhanced, breakable, colback=blue!3, colframe=blue, fonttitle=\bfseries,
  title={#1}, boxrule=0.6pt, left=6pt, right=6pt, top=4pt, bottom=4pt, arc=2pt}
\newtcolorbox{claim}[1]{enhanced, breakable, colback=white, colframe=black, fonttitle=\bfseries,
  title={#1}, boxrule=0.6pt, left=6pt, right=6pt, top=4pt, bottom=4pt, arc=2pt}
\newtcolorbox{keypoint}{enhanced, breakable, colback=green!6, colframe=green, boxrule=0.6pt,
  left=6pt, right=6pt, top=4pt, bottom=4pt, arc=2pt}
\newtcolorbox{caution}{enhanced, breakable, colback=red!5, colframe=red, boxrule=0.6pt,
  left=6pt, right=6pt, top=4pt, bottom=4pt, arc=2pt}
\newtcolorbox{task}[1]{enhanced, breakable, colback=yellow!12, colframe=yellow!60!black, fonttitle=\bfseries,
  title={#1}, boxrule=0.6pt, left=6pt, right=6pt, top=4pt, bottom=4pt, arc=2pt}

\lstset{
  basicstyle=\ttfamily\small, keywordstyle=\color{blue}\bfseries,
  commentstyle=\color{green}, stringstyle=\color{red},
  showstringspaces=false, breaklines=true, frame=single, rulecolor=\color{black!30},
  numbers=none, xleftmargin=6pt, framexleftmargin=6pt, language=Python
}
\newtheorem{theorem}{Theorem}
\newtheorem{lemma}{Lemma}
\newtheorem{proposition}{Proposition}
\theoremstyle{definition}
\newtheorem{definition}{Definition}
\newtheorem{exercise}{Exercise}

% =============================================================================
%  Notation — one convention for all lectures
%
%  scalars      x, y, w           plain italic
%  vectors      \vx \vy \vw       bold lower-case
%  matrices     \vX \vW \vSigma   bold upper-case
%  parameters   \vbeta \vtheta    bold Greek
%  transpose    \vX\T             \top
% =============================================================================
\newcommand{\R}{\mathbb{R}}
\newcommand{\E}{\mathbb{E}}
\newcommand{\Prob}{\mathbb{P}}
\DeclareMathOperator{\Var}{Var}
\DeclareMathOperator{\Cov}{Cov}
\DeclareMathOperator{\tr}{tr}
\DeclareMathOperator{\diag}{diag}
\DeclareMathOperator{\rank}{rank}
\DeclareMathOperator{\sign}{sign}
\DeclareMathOperator{\col}{col}
\DeclareMathOperator{\RSS}{RSS}
\DeclareMathOperator*{\argmin}{arg\,min}
\DeclareMathOperator*{\argmax}{arg\,max}
\newcommand{\norm}[1]{\left\lVert #1 \right\rVert}
\newcommand{\abs}[1]{\left\lvert #1 \right\rvert}
\newcommand{\inner}[2]{\left\langle #1,\, #2 \right\rangle}
\newcommand{\ind}[1]{\mathbf{1}\!\left\{#1\right\}}
\newcommand{\T}{^{\top}}
\newcommand{\inv}{^{-1}}
\newcommand{\half}{\tfrac{1}{2}}
\newcommand{\eps}{\varepsilon}
\newcommand{\Dset}{\mathcal{D}}
\newcommand{\Hset}{\mathcal{H}}
\newcommand{\Lcal}{\mathcal{L}}
\newcommand{\Ncal}{\mathcal{N}}

% bold vectors
\newcommand{\vx}{\mathbf{x}}   \newcommand{\vy}{\mathbf{y}}   \newcommand{\vz}{\mathbf{z}}
\newcommand{\vw}{\mathbf{w}}   \newcommand{\vu}{\mathbf{u}}   \newcommand{\vv}{\mathbf{v}}
\newcommand{\vb}{\mathbf{b}}   \newcommand{\vh}{\mathbf{h}}   \newcommand{\vr}{\mathbf{r}}
\newcommand{\vc}{\mathbf{c}}   \newcommand{\vf}{\mathbf{f}}   \newcommand{\vg}{\mathbf{g}}
\newcommand{\vp}{\mathbf{p}}   \newcommand{\vq}{\mathbf{q}}   \newcommand{\vs}{\mathbf{s}}
\newcommand{\va}{\mathbf{a}}   \newcommand{\vd}{\mathbf{d}}   \newcommand{\vt}{\mathbf{t}}
\newcommand{\vk}{\mathbf{k}}   \newcommand{\vm}{\mathbf{m}}   \newcommand{\vn}{\mathbf{n}}
\newcommand{\vzero}{\mathbf{0}} \newcommand{\vone}{\mathbf{1}}
% bold matrices
\newcommand{\vX}{\mathbf{X}}   \newcommand{\vY}{\mathbf{Y}}   \newcommand{\vZ}{\mathbf{Z}}
\newcommand{\vW}{\mathbf{W}}   \newcommand{\vH}{\mathbf{H}}   \newcommand{\vI}{\mathbf{I}}
\newcommand{\vA}{\mathbf{A}}   \newcommand{\vB}{\mathbf{B}}   \newcommand{\vC}{\mathbf{C}}
\newcommand{\vD}{\mathbf{D}}   \newcommand{\vP}{\mathbf{P}}   \newcommand{\vQ}{\mathbf{Q}}
\newcommand{\vU}{\mathbf{U}}   \newcommand{\vV}{\mathbf{V}}   \newcommand{\vS}{\mathbf{S}}
\newcommand{\vM}{\mathbf{M}}   \newcommand{\vK}{\mathbf{K}}   \newcommand{\vG}{\mathbf{G}}
% bold Greek
\newcommand{\vbeta}{\boldsymbol{\beta}}     \newcommand{\valpha}{\boldsymbol{\alpha}}
\newcommand{\vtheta}{\boldsymbol{\theta}}   \newcommand{\vmu}{\boldsymbol{\mu}}
\newcommand{\vSigma}{\boldsymbol{\Sigma}}   \newcommand{\vLambda}{\boldsymbol{\Lambda}}
\newcommand{\veps}{\boldsymbol{\varepsilon}} \newcommand{\vphi}{\boldsymbol{\phi}}
\newcommand{\vgamma}{\boldsymbol{\gamma}}   \newcommand{\vlambda}{\boldsymbol{\lambda}}
\newcommand{\vdelta}{\boldsymbol{\delta}}   \newcommand{\vnu}{\boldsymbol{\nu}}
\newcommand{\vxi}{\boldsymbol{\xi}}         \newcommand{\vpi}{\boldsymbol{\pi}}
\newcommand{\vomega}{\boldsymbol{\omega}}




\title{\textcolor{blue}{Lab 3 --- Trees, Random Forests and Gradient Boosting}\\[4pt]
\large Machine Learning I \,\textperiodcentered\, Tutorial}
\author{Rodrigo Grijalba (teaching assistant) \,\textperiodcentered\, Alexander Quispe (instructor)}
\date{Wednesday, September 23, 2026 \,\textperiodcentered\, 19:00--22:00}

\begin{document}
\maketitle

\begin{keypoint}
\textbf{Goal of the session.} Compute an information gain by hand and confirm scikit-learn picks the
same split; watch a single tree overfit and a bagged ensemble not; check the variance formula
$\rho\sigma^2+\frac{1-\rho}{B}\sigma^2$ on real trees; reproduce the AdaBoost trace of Lecture 5 to
three decimals; and tune the number of boosting rounds the only way that is honest.
\end{keypoint}

\section*{Materials}
\begin{itemize}[nosep]
  \item \texttt{lab\_03\_student.ipynb} / \texttt{lab\_03\_solution.ipynb}.
  \item Sources: CS 155 Problem Set 3 (decision trees, information gain, \texttt{min\_samples\_leaf}).
  \item Data: \texttt{load\_breast\_cancer} ($N=569$, $p=30$), bundled with scikit-learn.
\end{itemize}

\section{Impurity and Information Gain (35 min)}

\begin{task}{Task 1 --- three impurities, two lecture examples}
Write \texttt{entropy(p)}, \texttt{gini(p)}, \texttt{misclass(p)} for a vector of class proportions, and
\texttt{gain(parent, children, impurity)} where \texttt{children} is a list of class-count vectors.
Reproduce Lecture 4:
\begin{enumerate}[nosep]
  \item Parent $(8,4)$ split into $(6,1)$ and $(2,3)$: gains $0.1685$ (entropy), $0.1016$ (Gini), $0.0833$ (misclassification).
  \item Parent $(400,400)$: split A into $(300,100),(100,300)$ vs.\ split B into $(200,400),(200,0)$.
        Misclassification ties at $0.25$ after either split; Gini gives $0.375$ vs.\ $0.333$; entropy $0.811$ vs.\ $0.689$.
\end{enumerate}
\end{task}

\begin{task}{Task 2 --- the root split, by hand}
Write \texttt{best\_split(X, y, impurity)}: for every feature $j$ and every midpoint $s$ between consecutive
sorted values of $x_{\cdot j}$, compute the gain of the split $x_j\le s$; return the best $(j,s,\Delta)$.
Run it on the breast-cancer training set with Gini. Then fit \texttt{DecisionTreeClassifier(max\_depth=1)}
and read \texttt{tree\_.feature[0]} and \texttt{tree\_.threshold[0]}. They must agree.
\end{task}

\section{One Tree (25 min)}

\begin{task}{Task 3 --- overfitting, in two directions}
With an 80/20 split:
\begin{enumerate}[nosep]
  \item For \texttt{max\_depth} $\in\{1,2,\dots,12\}$, plot training and test error. Where does the test curve turn?
  \item For \texttt{min\_samples\_leaf} $\in\{1,2,5,10,20,50\}$ at unlimited depth, same plot. Explain, using
        Lecture 4's remark that a leaf's prediction has variance at most $\sigma^2/n_{\min}$, why this knob
        works.
  \item Grow the unrestricted tree and report: number of leaves, training error, test error. This is the
        low-bias, high-variance object that bagging is built for.
\end{enumerate}
\end{task}

\section{Bagging and the Variance Formula (45 min)}

\begin{task}{Task 4 --- bagging from scratch}
Implement \texttt{bagged\_trees(X, y, B, seed)}: $B$ bootstrap samples (\texttt{rng.integers(0, N, N)}),
one unrestricted tree each; return the list of trees and the list of index arrays. Predict by majority vote.
\begin{enumerate}[nosep]
  \item Plot test error against $B\in\{1,2,5,10,25,50,100,200\}$. Overlay the single-tree error from Task 3.
  \item Out-of-bag error: for each training point, vote only among trees whose bootstrap sample did not
        contain it. Compare with the test error and with \texttt{BaggingClassifier(oob\_score=True)}.
  \item Check the $e^{-1}$ claim: the mean fraction of distinct training points per bootstrap sample should be
        $\approx0.632$.
\end{enumerate}
\end{task}

\begin{task}{Task 5 --- measure $\rho$ and test the formula}
Treat each tree's prediction on the test set as a $\pm1$ vector. Compute the average pairwise correlation
$\hat\rho$ across the $B=100$ trees and the average per-tree variance $\hat\sigma^2$ (across test points).
Then compute the variance of the \emph{averaged} vote directly, and compare with
$\hat\rho\hat\sigma^2+\frac{1-\hat\rho}{B}\hat\sigma^2$. Report $\hat\rho$: it is the floor bagging cannot go below.
\end{task}

\begin{task}{Task 6 --- random forests lower $\rho$}
Repeat Task 5 with \texttt{RandomForestClassifier(max\_features=m)} for $m\in\{1,3,5\ (\approx\sqrt{30}),10,30\}$.
Tabulate $\hat\rho$, single-tree test error, and forest test error against $m$. The pattern to see: $\hat\rho$
falls as $m$ falls; single trees get worse; the forest gets better --- until $m$ is so small that the trees are
useless. Finish with permutation importance (\texttt{sklearn.inspection.permutation\_importance}) on the best forest:
which five features matter?
\end{task}

\section{Boosting (50 min)}

\begin{task}{Task 7 --- AdaBoost by hand}
Implement AdaBoost with decision stumps on a 1-D dataset, base learners $h(x)=\pm\sign(x-t)$ with $t$ at
midpoints. Run it on $x=(1,2,3,4,5)$, $y=(+1,+1,-1,-1,+1)$ for three rounds and print, per round, the
weights, the chosen stump, $\varepsilon_m$, $\alpha_m$, $Z_m$. Your numbers must match Lecture 5's table:
$\varepsilon=(0.2,0.25,0.333)$, $\alpha=(0.693,0.549,0.347)$, $Z=(0.8,0.866,0.943)$, and $\prod Z_m=0.653$.
Then verify the two identities of the proof: $Z_m=2\sqrt{\varepsilon_m(1-\varepsilon_m)}$, and
$\frac1N\sum_ie^{-y_if_3(x_i)}=\prod_mZ_m$.
\end{task}

\begin{task}{Task 8 --- early stopping on $M$}
On the breast-cancer data, fit \texttt{GradientBoostingClassifier(n\_estimators=2000, learning\_rate=$\nu$,
max\_depth=3)} for $\nu\in\{1,0.1,0.01\}$. Use \texttt{staged\_predict\_proba} to get the validation log loss after
every round on a held-out 20\% of the training set. Plot the three curves. For each $\nu$, report the $M$ that
minimizes validation loss and the test error of the model truncated there. Which $\nu$ wins, and at what cost in $M$?
\end{task}

\begin{task}{Task 9 --- three ensembles, one table}
Tune each of bagging, random forest and gradient boosting by 5-fold CV on the training set (a small grid
each --- $B$, $m$, and $(\nu,M,\text{depth})$ respectively). Report test accuracy, AUC and wall-clock fit time
for the three tuned models and for the single tuned tree. Two sentences on which you would deploy and why.
\end{task}

\section{Exercises (to submit before Lab 4)}
\begin{exercise}
Gradient boosting for regression, from scratch: with squared loss and stumps, implement the loop of
Lecture 5's algorithm (pseudo-residuals, fit, line search, shrinkage). Check that with $\nu=1$ it matches
\texttt{GradientBoostingRegressor(max\_depth=1, learning\_rate=1)} on the diabetes data after 50 rounds.
\end{exercise}
\begin{exercise}
Verify XGBoost's leaf formula in the simplest case: for squared loss, $g_i=f(\vx_i)-y_i$ and $k_i=1$, so
with $\lambda=0$ the optimal leaf value $-G_j/K_j$ is the mean residual in the leaf. Show this numerically
on one boosting round, then repeat with $\lambda=10$ and explain what changed.
\end{exercise}
\begin{exercise}
Prove that in AdaBoost the weighted error of $h_m$ under the \emph{updated} weights $\vw^{(m+1)}$ is exactly
$\tfrac12$. (Hint: split $Z_m$ into the correct and incorrect parts and use the value of $\alpha_m$.)
Verify it in your Task 7 code. What does this say about why the same stump is never chosen twice in a row?
\end{exercise}

\section*{Submission}
One executed notebook, \texttt{lab03\_lastname\_firstname.ipynb}, to the course drive before
Wednesday, September 30, 18:00.

\end{document}
